\documentclass[10pt,a4paper]{amsart}
\usepackage{amsmath,amssymb,amsthm,graphicx}
\usepackage{verbatim}
\newenvironment{reference}{\comment}{\endcomment}
%\newenvironment{reference}{}{}
\newenvironment{slogan}{\comment}{\endcomment}
\newenvironment{history}{\comment}{\endcomment}

% For commutative diagrams we use Xy-pic
\usepackage[all]{xy}

% We use 2cell for 2-commutative diagrams.
\xyoption{2cell}
% Theorem environments.
%
\theoremstyle{plain}
\newtheorem{theorem}[subsection]{Theorem}
\newtheorem{proposition}[subsection]{Proposition}
\newtheorem{lemma}[subsection]{Lemma}

\theoremstyle{definition}
\newtheorem{definition}[subsection]{Definition}
\newtheorem{example}[subsection]{Example}
\newtheorem{exercise}[subsection]{Exercise}
\newtheorem{situation}[subsection]{Situation}

\theoremstyle{remark}
\newtheorem{remark}[subsection]{Remark}
\newtheorem{remarks}[subsection]{Remarks}

\numberwithin{equation}{subsection}

% Macros
%
\def\lim{\mathop{\mathrm{lim}}\nolimits}
\def\colim{\mathop{\mathrm{colim}}\nolimits}
\def\Spec{\mathop{\mathrm{Spec}}}
\def\Hom{\mathop{\mathrm{Hom}}\nolimits}
\def\Ext{\mathop{\mathrm{Ext}}\nolimits}
\def\SheafHom{\mathop{\mathcal{H}\!\mathit{om}}\nolimits}
\def\SheafExt{\mathop{\mathcal{E}\!\mathit{xt}}\nolimits}
\def\Sch{\mathit{Sch}}
\def\Mor{\mathop{\mathrm{Mor}}\nolimits}
\def\Ob{\mathop{\mathrm{Ob}}\nolimits}
\def\Sh{\mathop{\mathit{Sh}}\nolimits}
\def\NL{\mathop{N\!L}\nolimits}
\def\CH{\mathop{\mathrm{CH}}\nolimits}
\def\proetale{{pro\text{-}\acute{e}tale}}
\def\etale{{\acute{e}tale}}
\def\QCoh{\mathit{QCoh}}
\def\Ker{\mathop{\mathrm{Ker}}}
\def\Im{\mathop{\mathrm{Im}}}
\def\Coker{\mathop{\mathrm{Coker}}}
\def\Coim{\mathop{\mathrm{Coim}}}

% Boxtimes
%
\DeclareMathSymbol{\boxtimes}{\mathbin}{AMSa}{"02}

%
% Macros for moduli stacks/spaces
%
\def\QCohstack{\mathcal{QC}\!\mathit{oh}}
\def\Cohstack{\mathcal{C}\!\mathit{oh}}
\def\Spacesstack{\mathcal{S}\!\mathit{paces}}
\def\Quotfunctor{\mathrm{Quot}}
\def\Hilbfunctor{\mathrm{Hilb}}
\def\Curvesstack{\mathcal{C}\!\mathit{urves}}
\def\Polarizedstack{\mathcal{P}\!\mathit{olarized}}
\def\Complexesstack{\mathcal{C}\!\mathit{omplexes}}
% \Pic is the operator that assigns to X its picard group, usage \Pic(X)
% \Picardstack_{X/B} denotes the Picard stack of X over B
% \Picardfunctor_{X/B} denotes the Picard functor of X over B
\def\Pic{\mathop{\mathrm{Pic}}\nolimits}
\def\Picardstack{\mathcal{P}\!\mathit{ic}}
\def\Picardfunctor{\mathrm{Pic}}
\def\Deformationcategory{\mathcal{D}\!\mathit{ef}}
\usepackage{cleveref}
\crefname{originallemma}{original lemma}{original lemmas}
\Crefname{originallemma}{Original lemma}{Original lemmas}
\makeatletter
\newlabel{lemma-compose-finite-type}{{00F4}{source}{}{sourcepointer-1}{}}
\newlabel{lemma-compose-finite-type@cref}{{[originallemma][0][]00F4}{[1][0][]source}{}{}{}}
\makeatother
\crefname{originallemma}{original lemma}{original lemmas}
\Crefname{originallemma}{Original lemma}{Original lemmas}
\makeatletter
\newlabel{lemma-local-isomorphism}{{00QR}{source}{}{sourcepointer-2}{}}
\newlabel{lemma-local-isomorphism@cref}{{[originallemma][0][]00QR}{[1][0][]source}{}{}{}}
\makeatother
\title{An isomorphism of localisations of finitely presented algebras spreads out to an isomorphism after inverting one element}
\author{The Stacks Project Authors}
\date{stacks-project:algebra.tex:local-snapshot-20261012; GNU FDL 1.2 or later}
\begin{document}
\maketitle
\noindent\textit{Editorial note.} All delivered mathematical text is the author's own, copied from the pinned Stacks Project chapter file algebra.tex; no statement or proof step was rewritten, reordered or removed.\par\medskip
\noindent\textit{Editorial note.} Preamble provenance, as required for this source: the frozen snapshot stages the chapter's own preamble file beside it (preamble.tex, 217 lines, 8116 bytes, the same shared Stacks-wide preamble the published chow and etale-cohomology cards use); the first line of the pinned chapter file is \textbackslash{}input\{preamble\}, so no preamble text lives inside the chapter itself. The card takes that file's lines 8-18 (the verbatim package and the three \textbackslash{}comment-based environments reference/slogan/history at lines 9-12, section 13.2, together with the Xy-pic setup) and its lines 154 to the runtime end of the file (the native theorem declarations and mathematical macros). Not one character was changed, no package was added, no macro was defined, and the excerpt serves only as preamble, never as a span coordinate.\par\medskip
\noindent\textit{Editorial note.} The same-chapter pointers are located at their \textbackslash{}label lines inside the pinned chapter file; the cross-chapter pointers are located in the project's own tag list staged in the snapshot (sources/stacks-chow/tags.txt), at the verbatim lines carrying their canonical tags.\par\medskip
\noindent\textit{Editorial note.} The counted proof refers to two same-chapter statements by label; both are recorded as navigation stubs pointed at their \textbackslash{}label lines, and the arguments consuming them stay inside the counted span.\par\medskip
\noindent\textit{Editorial note.} The counted spans contain no \textbackslash{}cite, so the card carries no bibliography.\par\medskip
\noindent\textit{Allowed standard prerequisites (editorial).}\begin{itemize}
\item Clearing denominators in an isomorphism of localisations. Used at the start of the counted proof (lines 31852-31860). Produces elements outside the primes such that the two presentations become isomorphic after inverting them.
\item Splitting off a factor from an isomorphism of localisations. Invoked at the end of the counted proof (lines 31874-31876). Reduces the general case to a product decomposition, from which a localised isomorphism is extracted.
\end{itemize}\medskip
\ifdefined\hypertarget\hypertarget{sourcepointer-1}{}\fi
\noindent\textit{Source reference.} Original lemma 00F4: sources/stacks-algebra/pinned-original-source.tex, source lines 757--757. Navigation pointer for \textbackslash{}ref\{lemma-compose-finite-type\} used inside the counted span (line 31872); the argument that consumes it is retained in the span.\par\medskip
\ifdefined\hypertarget\hypertarget{sourcepointer-2}{}\fi
\noindent\textit{Source reference.} Original lemma 00QR: sources/stacks-algebra/pinned-original-source.tex, source lines 31798--31798. Navigation pointer for \textbackslash{}ref\{lemma-local-isomorphism\} used inside the counted span (line 31874); the argument that consumes it is retained in the span.\par\medskip
\begin{lemma}
\label{lemma-isomorphic-local-rings}
Let $R$ be a ring.
Let $S$, $S'$ be of finite presentation over $R$.
Let $\mathfrak q \subset S$ and $\mathfrak q' \subset S'$
be primes. If $S_{\mathfrak q} \cong S'_{\mathfrak q'}$ as
$R$-algebras, then there exist $g \in S$, $g \not \in \mathfrak q$
and $g' \in S'$, $g' \not \in \mathfrak q'$ such that
$S_g \cong S'_{g'}$ as $R$-algebras.
\end{lemma}

\begin{proof}
Let $\psi : S_{\mathfrak q} \to S'_{\mathfrak q'}$ be the isomorphism
of the hypothesis of the lemma.
Write $S = R[x_1, \ldots, x_n]/(f_1, \ldots, f_r)$ and
$S' = R[y_1, \ldots, y_m]/J$.
For each $i = 1, \ldots, n$ choose a fraction
$h_i/g_i$ with $h_i, g_i \in R[y_1, \ldots, y_m]$
and $g_i \bmod J$ not in $\mathfrak q'$ which represents
the image of $x_i$ under $\psi$. After replacing
$S'$ by $S'_{g_1 \ldots g_n}$ and
$R[y_1, \ldots, y_m, y_{m + 1}]$ (mapping $y_{m + 1}$ to $1/(g_1\ldots g_n)$)
we may assume that $\psi(x_i)$ is the image of some
$h_i \in R[y_1, \ldots, y_m]$. Consider the elements
$f_j(h_1, \ldots, h_n) \in R[y_1, \ldots, y_m]$.
Since $\psi$ kills each $f_j$ we see that
there exists a $g \in R[y_1, \ldots, y_m]$, $g \bmod J \not \in \mathfrak q'$
such that $g f_j(h_1, \ldots, h_n) \in J$ for each $j = 1, \ldots, r$.
After replacing $S'$ by $S'_g$ and
$R[y_1, \ldots, y_m, y_{m + 1}]$ as before we may assume that
$f_j(h_1, \ldots, h_n) \in J$. Thus we obtain a ring map
$S \to S'$, $x_i \mapsto h_i$ which induces $\psi$ on local rings.
By Lemma \ref{lemma-compose-finite-type}
the map $S \to S'$ is of finite presentation.
By Lemma \ref{lemma-local-isomorphism}
we may assume that $S' = S \times C$. Thus localizing $S'$ at the
idempotent corresponding to the factor $C$ we obtain the result.
\end{proof}

\end{document}
